\section{Quantum two-term Weyl law at fixed magnetic field}
\label{sec:quantum-weyl}

The identity in \cref{thm:classical-clock} is a phase-volume statement, not
a quantum asymptotic.  We now show that it survives quantization at the two
growing orders of the Riemann--von Mangoldt clock.  Throughout this section,
the H\'enon parameter \(a\), the iterate \(n\), and the magnetic field \(B\)
are fixed while \(E\to\infty\).

\subsection{The confining magnetic operator}

On \(C_c^\infty(\R^2)\), consider
\begin{equation}
 \mathfrak h_{a,n,B}[u]
 =\frac12\|(-\ii\nabla-A_B)u\|_2^2
  +\int_{\R^2}V_{a,n}(q)|u(q)|^2\dd q .
 \label{eq:magnetic-form}
\end{equation}
The form is closable; we use its closure and denote the corresponding
Friedrichs operator by
\(\HH_{a,n,B}\).

\begin{proposition}[Self-adjointness and discreteness]
\label{prop:compact-resolvent}
For every fixed \(a>-1\), \(n\geq1\), and \(B\in\R\), the form in
\cref{eq:magnetic-form} is densely defined, closable, and lower
semibounded.  Its closure is a closed form, and the associated Friedrichs
realization is self-adjoint and has compact resolvent.
\end{proposition}

\begin{proof}[Proof sketch]
The polynomial automorphism \(\Psi_{a,n}\) is proper, so
\(V_{a,n}(q)\to\infty\) as \(|q|\to\infty\).  On a fixed compact set the
linear vector potential \(A_B\) is bounded; hence a form-norm bounded family
is bounded in the ordinary local \(H^1\) norm.  Rellich compactness controls
the family on that compact set, whereas
\begin{equation}
 \int_{\{V_{a,n}>M\}}|u|^2
 \leq M^{-1}\int_{\R^2}V_{a,n}|u|^2
 \label{eq:potential-tail}
\end{equation}
controls the tail uniformly as \(M\to\infty\).  The full argument is given in
\cref{app:compactness}; see also the standard magnetic-form framework in
\citep{Kato1995,LeinfelderSimader1981,BravermanMilatovicShubin2002,MazyaShubin2005}.
\end{proof}

Let
\begin{equation}
 N_{a,n,B}(E)=\#\{k:\lambda_k(\HH_{a,n,B})\leq E\},
 \label{eq:quantum-count-definition}
\end{equation}
with multiplicity.

\begin{theorem}[Fixed-iterate magnetic Weyl law]
\label{thm:magnetic-weyl}
Let \(a>-1\), \(a\neq0\), and let \(n\geq1\) and \(B\in\R\) be fixed.  Put
\(D=2^n\).  Then
\begin{equation}
 \boxed{
 N_{a,n,B}(E)
 =\frac{E}{2\pi}\log\frac{E}{2\pi}
  -\frac{E}{2\pi}
  +\order{E^{3/4}(\log E)^{1+D/2}} }
 \label{eq:quantum-weyl}
\end{equation}
as \(E\to\infty\).  The implied constant may depend on \(a,n,B\).  In
particular,
\begin{equation}
 N_{a,n,B}(E)
 =\frac{E}{2\pi}\log\frac{E}{2\pi}
  -\frac{E}{2\pi}+o(E).
 \label{eq:two-growing-terms}
\end{equation}
\end{theorem}

For the radial control \(a=0\), the map \(\Psi_{0,n}\) is orthogonal and the
same argument applies with effective degree one; for example, it gives the
safe bound
\begin{equation}
 N_{0,n,B}(E)
 =\frac{E}{2\pi}\log\frac{E}{2\pi}
  -\frac{E}{2\pi}
  +\order{E^{3/4}(\log E)^{3/2}}.
 \label{eq:radial-magnetic-weyl}
\end{equation}
No optimality is asserted for either remainder.

\subsection{Proof architecture}

We record the estimates that determine the stated error and defer their
proofs to \cref{app:poly-geometry,app:magnetic-count,app:weyl-assembly}.
Write
\begin{equation}
 L=\log(E/2\pi),\qquad R=\sqrt{L/\pi},\qquad
 \Omega_E=\{q:|\Psi_{a,n}(q)|<R\}.
 \label{eq:allowed-domain}
\end{equation}
Area preservation gives \(|\Omega_E|=L\).  Because both
\(\Psi_{a,n}\) and its inverse are degree-\(D\) polynomial maps,
\begin{equation}
 \operatorname{length}(\partial\Omega_E)
 =O_{a,n}(L^{D/2}),
 \qquad
 \sup_{\Psi_{a,n}^{-1}(B_{R+1})}|\nabla\Phi_{a,n}|
 =O_{a,n}(L^{D/2}),
 \label{eq:allowed-geometry}
\end{equation}
where \(\Phi_{a,n}=\pi|\Psi_{a,n}|^2\).

Tile \(\R^2\) by squares \(Q\) of side
\begin{equation}
 \ell=E^{-1/4}.
 \label{eq:square-scale}
\end{equation}
For strict local counts, with the count set equal to zero at nonpositive
energy, Dirichlet--Neumann bracketing gives
\begin{equation}
 \sum_Q n_{B,D,Q}(E-V_Q^+)
 \leq N_{<}(E)
 \leq\sum_Q n_{B,N,Q}(E-V_Q^-).
 \label{eq:dn-bracketing}
\end{equation}
Here \(V_Q^-\) and \(V_Q^+\) are the extrema of \(V_{a,n}\) on \(Q\), and
the Neumann operator uses the covariant boundary condition
\(\nu\cdot(-\ii\nabla-A_B)u=0\).

The magnetic field does not change the leading local count.  On a square
centered at \(c_Q\), a local gauge removes \(A_B(c_Q)\), leaving a residual
of norm \(O_B(\ell)\).  Min--max comparison with the free square therefore
gives, uniformly for \(0\leq A\leq E\),
\begin{equation}
 n_{B,D/N,Q}(A)
 =\frac{\ell^2 A}{2\pi}+O_B(\ell\sqrt E+1).
 \label{eq:local-magnetic-count}
\end{equation}
The number of squares entering the upper sum is
\begin{equation}
 M(E)=O_{a,n}\!\left(
 \frac{L}{\ell^2}+\frac{L^{D/2}}{\ell}+1\right).
 \label{eq:relevant-squares}
\end{equation}
Consequently, the accumulated local lattice error is contained in
\(O_{a,n,B}(E^{3/4}L^{1+D/2})\).

The remaining issue is the potential variation on a boundary square.  A
first-exit bootstrap shows that every square meeting \(\Omega_E\) is mapped
into \(B_{R+1}\) for sufficiently large \(E\).  The second estimate in
\cref{eq:allowed-geometry} then yields
\begin{equation}
 V_Q^+-V_Q^-=O_{a,n}(E\ell L^{D/2}).
 \label{eq:potential-oscillation}
\end{equation}
The relevant-square union has area \(O_{a,n}(L)\), so the gap between the
upper and lower principal Riemann sums is
\begin{equation}
 O_{a,n}(E\ell L^{1+D/2})
 =O_{a,n}(E^{3/4}L^{1+D/2}).
 \label{eq:riemann-sum-error}
\end{equation}
The common principal integral is precisely the classical quantity in
\cref{thm:classical-clock}.  Finally,
\(N_{<}(E)\leq N_{\leq}(E)\leq N_{<}(E+1)\); the change of the comparison
function over this unit interval is \(O(\log E)\).  This proves
\cref{thm:magnetic-weyl}.  General Weyl results for confining and magnetic
Schr\"odinger operators provide context, but the explicit remainder here
comes from the preceding bracketing estimates
\citep{Rozenblum1974,HelfferRobert1982,Tachizawa1992,Ivrii2016}.

\begin{remark}[Exactly what the theorem resolves]
\label{rem:mean-resolution}
Equation \eqref{eq:quantum-weyl} resolves the coefficients of \(E\log E\)
and \(E\).  The exact classical comparison function in
\cref{eq:exact-classical-clock} contains \(+1\), but it is intentionally not
displayed as a quantum term because the quantum error is much larger than a
constant.  The theorem therefore determines neither the Riemann counting
constant \(7/8\) nor an analogue of the oscillatory term \(S(E)\)
\citep{TitchmarshHeathBrown1986}.  It also gives no uniform result when
\(a=a(E)\), \(n=n(E)\), or \(B=B(E)\).
\end{remark}

Identical growing counting terms do not imply identical spectra.  The next
section proves this distinction for the centered one-step scalar pair: its
ground state is strictly displaced and its relative heat trace has a nonzero
uniform asymptotic.
